\section{Notation}

\begin{list}{}{}
    \setlength\itemsep{0em}
    \item $O(f(n))$: Upper bound - Rate of growth for worst-case input of size $n$
    \item $\Omega(f(n))$: Lower bound - Rate of growth for best-case input of size $n$
    \item $\Theta(f(n))$: Tight bound - When $O(f(n)) = \Omega(f(n))$
    \item $o(g(n))$: Strictly upper bound - $f(n)$ grows slower than $g(n)$
    \item $\omega(g(n))$: Strictly lower bound - $f(n)$ grows faster than $g(n)$ 
\end{list}

When finding complexity, only the fastest growing term is kept, and constants are dropped.
Only $\Theta$ needs be tight, the others just have to be upper/lower bounds, even if they are far from tight.


\subsection{Hierarchy of growth rates}

\[
    1 < \lg \lg n < \lg n < \lg^k n < n^{k \in ]0,1[} < n < n \lg n < n^2 < n^3 < 2^n < n! < n^n
\]

Note: $\sqrt{n} = n^{1/2}$

\subsection{Ackermann functions}
\begin{list}{}{}
    \setlength\itemsep{0em}
    \item $A_m(n)$: The Ackermann function is practically infinite for $n > 4$.
    \item $\alpha(n)$: The Inverse Ackermann function is $\leq 4$ for all practical values $n$.
\end{list}

\subsubsection{Use}

Example of the Ackermann function:

Given that $A_2(1)=7$:
\[
    A_3(1) = A_2(A_2(1)) = A_2(7) = 2^{7+1}-1=2047
\]

This means that:
\[
    A_4(1) = A_3(A_3(1)) = A_3(2047) = A_2(A_2(...A_2(1)...)) \text{ [2047 times]} \approx \infty
\]
